\documentclass{article}
\usepackage{iclr2027_conference,times}

\usepackage{amsmath,amssymb,amsthm}
\usepackage{booktabs}
\usepackage{multirow}
\usepackage{graphicx}
\usepackage{microtype}
\usepackage{xcolor}
\usepackage{tikz}
\usetikzlibrary{arrows.meta,positioning,decorations.pathmorphing}
\usepackage[colorlinks=true,linkcolor=blue,citecolor=blue]{hyperref}
\usepackage{url}

% -------------------------------- macros --
\newcommand{\E}{\mathbb{E}}
\newcommand{\ind}[1]{\mathbf{1}\!\left[#1\right]}
\newcommand{\sg}{\mathrm{sg}}
\newcommand{\cube}{c}
\newcommand{\Reach}{\mathrm{reach}}
\newcommand{\Visit}{\mathrm{visit}}
\newcommand{\Nhat}{\hat N}

\newtheorem{proposition}{Proposition}


\title{Region-Based Successor Measures for Offline Robot Manipulation}
\author{Anonymous}

% Uncomment for the camera-ready version.
% \iclrfinalcopy

\begin{document}
\maketitle

% ============================================================================
\section{Successor densities versus region occupancy}
\label{sec:bounded}

\subsection{Successor measures and forward-backward representations}
\label{sec:prelim}

\begin{figure}[htb]
\centering
\includegraphics[width=0.98\textwidth]{figures/0_playdata_strip.png}
\caption{The data behind the failure: one trajectory from \texttt{cube-single-play-v0}, rendered by replaying the dataset's stored simulator states.
The collector moves the cube, pauses, re-grasps, and returns to the same configurations repeatedly; between manipulations the cube does not move at all.
The cube is stationary on $48\%$ of all transitions, so the discounted occupancy a successor measure must predict concentrates at these resting configurations (\S\ref{sec:playdata}).}
\label{fig:strip}
\end{figure}

\paragraph{Goal-conditioned RL.}
We consider a Markov decision process with state space $\mathcal{S}$, action space $\mathcal{A}$, transition kernel $P$, and discount $\gamma\in(0,1)$.
In goal-conditioned RL, a goal $g$ defines the sparse reward $r_g(s')=\ind{s'\in S_g}$ for some goal set $S_g$, and the standard reachability value of a policy $\pi$ (the goal-conditioned value under the absorbing convention) is $\E_\pi[\gamma^{T_g}]$, where $T_g$ is the hitting time of $S_g$ \citep{kaelbling1993learning,andrychowicz2017hindsight}.
Offline, the agent learns from a fixed dataset of transitions; hindsight relabeling supplies goals from observed states.

\paragraph{Successor measures.}
The successor measure of $\pi$ is its discounted state occupancy, viewed as a measure over sets $X\subseteq\mathcal{S}$:
\begin{equation}
\label{eq:sm}
M^{\pi}(s,a,X) = \E_{\pi}\!\Big[ \sum_{t\ge 0}\gamma^{t}\ind{s_{t+1}\in X} \,\Big|\, s_{0}=s,\,a_{0}=a \Big].
\end{equation}
Intuitively, $M^\pi(s,a,X)$ counts the discounted time the policy will spend in $X$.
Its usefulness is that any reward's value is an integral against it, $Q_r^{\pi}(s,a)=\int r(x)\,dM^{\pi}(s,a,x)$: dynamics are learned once, rewards are supplied later \citep{dayan1993improving,blier2021learning}.

Its Bellman equation, written over sets, is
\begin{equation}
\label{eq:density-bellman}
M^{\pi_z}(s,a,X) = P(s'\in X\mid s,a)+\gamma\,\E_{s',a'\sim\pi_z}\big[M^{\pi_z}(s',a',X)\big].
\end{equation}

\paragraph{Forward-backward representations and zero-shot RL.}
FB methods \citep{touati2021learning,touati2023does} factorize the successor measure's density with respect to a reference distribution $\rho$ (in practice, the dataset marginal):
\begin{equation}
\label{eq:fb}
M^{\pi_z}(s,a,X) \approx \int_X F(s,a,z)^{\top}\psi(x)\,d\rho(x).
\end{equation}
Here $z$ is a task vector indexing a family of policies $\pi_z$, and $x$ is a query state, independent of $z$.
At test time, a reward $r$ observed on dataset samples is converted to a task vector by least-squares reward inference, $z_r \propto \E_{x\sim\rho}[r(x)\psi(x)]$, and the agent acts greedily on $Q_r(s,a)\approx F(s,a,z_r)^{\top}z_r$ with no further training.
This is the zero-shot RL protocol we evaluate throughout.

\subsection{Why pointwise density estimation is problematic}
\label{sec:failure}

\begin{proposition}[Point densities are singular; region targets are bounded]
\label{prop:fixed-points}
Suppose the dynamics are deterministic, $s'=f(s,a)$, and the reference distribution $\rho$ is non-atomic.
Then the one-step transition measure $\delta_{f(s,a)}$ has no density with respect to $\rho$.
For any measurable region $G$, however, $P(G\mid s,a)=\ind{f(s,a)\in G}\in\{0,1\}$, and the region-occupancy Bellman operator obtained from Equation~\ref{eq:density-bellman} is a $\gamma$-contraction on bounded functions.
(Proof in Appendix~\ref{app:derivations}.)
\end{proposition}

The first half is the problem: FB regresses the density $m_z=dM^{\pi_z}/d\rho$, whose one-step Bellman target contains the density of $\delta_{f(s,a)}$ with respect to $\rho$, which does not exist.
The $\rho$-weighted squared error remains numerically defined, but it is minimized by densities that grow without bound at the observed transitions.
The singularity follows from deterministic dynamics alone; it is not created by manipulation data.
Manipulation play amplifies its practical effect (\S\ref{sec:playdata}): long stationary segments repeatedly present the same matched transition to the non-saturating term of FB's surrogate loss.

Standard FB works around the singular term with a surrogate loss \citep{touati2021learning},
\begin{equation}
\label{eq:fb-loss}
\mathcal{L}_{\mathrm{FB}} = \tfrac{1}{2}\E_{x\sim\rho}\big[(m_z-\gamma\bar m_z)^2\big] -\E\big[m_z(s,a,s')-\gamma\bar m_z(s,a,s')\big] +\lambda_{\mathrm{ortho}}\mathcal{L}_{\mathrm{ortho}}.
\end{equation}
This has the structure of a contrastive objective: random query states are negatives and receive a squared residual, while the observed successor is the positive pair and receives a term that rewards increasing its score \citep{eysenbach2022contrastive}.
But where a contrastive loss saturates its positive term (through a log-sigmoid), the term here is \emph{linear}: increasing the matched score always reduces the loss, at any scale, and on play data the same matched pairs recur for dozens of consecutive steps.
The only component of the objective that opposes the growth is the orthonormality penalty on $\psi$, a constraint on representation geometry acting as an implicit normalizer.

\subsection{Learning occupancy of finite regions}
\label{sec:regions}

The second half of Proposition~\ref{prop:fixed-points} is the method.
The environment provides a low-dimensional outcome map $\cube(x)$ (here the object position) and a radius $\epsilon$, which define regions $B_\epsilon(g)=\{x:\|\cube(x)-\cube(g)\|\le\epsilon\}$.
Define the indicator reward
\begin{equation}
\label{eq:hit}
h(s',g)=\ind{\|\cube(s')-\cube(g)\|\le\epsilon}\in\{0,1\},
\end{equation}
and regress the region occupancy $M(s,a,B_\epsilon(g))$ instead of the density $m(s,a,g)$.
Evaluated on the region $B_\epsilon(g)$, the immediate term of Equation~\ref{eq:density-bellman} is $P(s'\in B_\epsilon(g)\mid s,a)$, whose sample estimate is the indicator $h(s',g)$, a bounded, sparse, goal-conditioned reward of the kind offline GCRL trains on routinely.
The loss is full-matrix TD regression over the batch,
\begin{equation}
\label{eq:matrix-loss}
\mathcal{L}=\tfrac{1}{n^2}\sum\nolimits_{ij}\big(Q_{ij}-y_{ij}\big)^2, \qquad Q_{ij}=F(s_i,a_i,z_j)^{\top}\psi(g_j),
\end{equation}
with no linear term and no orthonormality penalty: the $0/1$ reward supplies the scale that the penalty previously had to impose.
We have replaced pointwise density regression with a bounded region-occupancy TD objective.

\paragraph{Two backups, one per component.}
A finite event admits two TD targets, distinguished by what happens after the first hit:
\begin{equation}
\label{eq:backups}
y^{\Reach}_{ij}=h_{ij}+\gamma(1-h_{ij})\,\bar Q_{ij}, \qquad\qquad y^{\Visit}_{ij}=h_{ij}+\gamma\,\bar M_{ij}.
\end{equation}
The first treats the region as absorbing, exactly as in standard goal-conditioned RL; the second does not, so occupancy keeps accumulating.

For a fixed continuation policy, both operators are $\gamma$-contractions with unique fixed points $Q^{\Reach}(s,a,g)=\E[\gamma^{T_g}\ind{T_g<\infty}]$ and $M^{\Visit}(s,a,g)=\E[\sum_{t\ge0}\gamma^t h(s_{t+1},g)]$, where $T_g=\min\{t\ge 0: s_{t+1}\in B_\epsilon(g)\}$ is the hitting time of $B_\epsilon(g)$ along the successor sequence (Appendix~\ref{app:derivations}).
These are the two components that \S\ref{sec:playdata} separates empirically: $Q^{\Reach}\in[0,1]$ is the standard reachability value, the small component that carries the control signal, and $M^{\Visit}\in[0,1/(1-\gamma)]$ is the successor measure of the region, whose magnitude on play data comes mostly from occupancy after the first arrival.
We refer to them below as the \textbf{reachability value} and the \textbf{region occupancy}.

\subsection{Reachability pretraining and occupancy fine-tuning}

The data statistics motivate training them in that order.
Phase 1 trains the reachability value with the policy's goal equal to the goal in the TD loss: a standard sparse-reward goal-conditioned RL problem in which every unit of value reflects reaching.
Phase 2 warm-starts from that checkpoint, freezes $\psi$, switches to the non-absorbing backup, and samples the task $z$ and the goal $g$ independently, $M(s,a,z,g)=F(s,a,z)^{\top}\psi(g)$, the separated form that reward integration requires (Equation~\ref{eq:fb}).
Phase 2 adds the post-arrival occupancy that reward integration needs on top of a goal representation that phase 1 has already validated for control.
The two-stage procedure is heuristic: occupancy does not factor exactly into reachability and recurrence for finite regions (Appendix~\ref{app:twostage}), so phase 2 optimizes occupancy directly with TD rather than composing estimated factors.

\paragraph{Freezing the goal embedding.}
Phase 2 trains the actor DDPG-style against the occupancy critic while the goal embedding $\psi$ receives no gradient.
The freeze preserves the goal representation validated in phase 1: \S\ref{sec:frozenrep} shows that training $\psi$ against concentrated targets corrupts it, and collapse under the density objective is twice as fast with $\psi$ trainable as with it frozen.
This resembles fixed-feature successor learning \citep{barreto2017successor} and slow-target stabilization in actor-critic methods \citep{lillicrap2016continuous}, although it is distinct from the ordinary stop-gradient applied to TD bootstrap targets; the connections are discussed in Appendix~\ref{app:freeze}.

\subsection{Test-time reward inference}
\label{sec:readout}

The readout follows from the same change of estimand.
Any reward's value is $Q_r^{\pi}(s,a)=\E_{g\sim\rho}[r(g)\,m^{\pi}(s,a,g)]$.
For regions small enough that the density is approximately constant on the ball, $M^{\pi}(s,a,B_\epsilon(g))\approx m^{\pi}(s,a,g)\,N_\epsilon(g)$ with $N_\epsilon(g)=\rho(B_\epsilon(g))$, so
\begin{equation}
\label{eq:readout}
Q_r^{\pi}(s,a) \approx \E_{g\sim\rho}\!\left[\frac{r(g)}{N_\epsilon(g)}\, M^{\pi}(s,a,B_\epsilon(g))\right], \qquad z_r=\E_{g\sim\rho}\!\left[\frac{r(g)}{N_\epsilon(g)}\,\psi(g)\right],
\end{equation}
where the second expression is the induced task embedding for the bilinear model and $N_\epsilon$ is estimated by a small trained head $\Nhat$ (\S\ref{sec:rq3}).
Equation~\ref{eq:readout} has one training-side consequence: inferred embeddings $z_r$ are averages of goal embeddings, so the distribution of $z$ seen in training must include such averages; we add reward-induced embeddings to the task-sampling distribution (\S\ref{sec:rq2}, Appendix~\ref{app:rewards}).

\section{Why the issue matters on manipulation play}
\label{sec:playdata}

What is different about manipulation play data?
In the domains where FB was developed, unsupervised-exploration locomotion and maze navigation, the variable that defines tasks is the agent's own configuration, and it changes at essentially every step: a walker's joints never stop, a maze agent keeps moving.
Occupancy under such data is diffuse.
Manipulation play data is the opposite.
The goal-relevant state variable is the \emph{object's} position, and the object only moves while grasped (Figure~\ref{fig:strip}).

\begin{figure}[t]
\centering
\includegraphics[width=0.95\textwidth]{figures/playdata.png}
\caption{Manipulation play data concentrates occupancy where the object rests.
\emph{Left:} cube traces of three play trajectories, the cube rests, is carried, and rests again.
\emph{Right:} discounted occupancy of a $4$\,cm goal ball accumulated \emph{after} its first arrival (within-trajectory Monte Carlo, $\gamma=0.99$): the mean is $30.6$ discounted visits, out of a maximum possible $1/(1-\gamma)=100$.
Most of the magnitude of the successor measure on this data accrues after the first arrival, not from reaching.}
\label{fig:playdata}
\end{figure}

Three statistics quantify this (Figure~\ref{fig:playdata}).
First, on \texttt{cube-single-play-v0} the cube position is bit-identical across $48\%$ of all transitions, with a median stationary run of $44$ consecutive steps.
Second, a $4$\,cm ball around a random dataset state contains only $1.3\%$ of the data, yet the discounted occupancy a trajectory accumulates in that ball \emph{after first arriving} averages $25$-$31$ visits out of a possible $100$: most of the successor measure's magnitude on this data comes from the object sitting inside the goal region, not from reaching it.
Third, post-arrival occupancy is predictable from the goal alone: a regressor given only the goal predicts Monte-Carlo occupancy at rank correlation $0.60$-$0.63$, while the reachability value alone predicts it at $0.06$.
This suggests that the component dominating the measure's magnitude contains little action-selection information, and that most of that information is in the smaller first-arrival component.

The structure strengthens with every additional object, because play manipulates one object at a time while every other object rests.
On \texttt{cube-double-play-v0}, each cube is stationary on $72\%$ of transitions and post-arrival occupancy of the goal region averages $32$; on \texttt{scene-play-v0}, the drawer and window are stationary on $91\%$ of transitions, the buttons on $99.6\%$, and post-arrival occupancy averages $37$, while the conjunctive goal region (all objects in place) covers only $0.03$-$0.4\%$ of the data (Appendix~\ref{app:playstats}).
If this mechanism is what breaks density estimation, methods built on occupancy densities should degrade most where the concentration is strongest, and they do: FB scores $0.42$ on \texttt{cube-single} but $0.05$ on \texttt{scene} (Table~\ref{tab:main}).

\subsection{A frozen-representation experiment isolates the objective}
\label{sec:frozenrep}

To rule out representation quality as an explanation, we freeze $\psi$ at a checkpoint that already supports $0.97$ task success and train only the density estimator on top of it.
Without the orthonormality term, the critic's scale grows from $154$ to $5.4\times10^4$ within $65$k steps and mean success decreases monotonically from $0.98$ to $0.42$ by $100$k steps, three seeds (Figure~\ref{fig:diverge}); letting $\psi$ train merely accelerates the collapse ($0.97\to0.07$ by $50$k).
With the orthonormality term, scale is stable but success reaches only $0.225$; a fully tuned FB implementation reaches $\approx0.42$, against $0.928$ for the architecture-matched GCIQL baseline.
The estimator, not the representation, is the bottleneck.

\begin{figure}[t]
\centering
\begin{minipage}[t]{0.48\textwidth}
\centering
\includegraphics[width=\textwidth]{figures/scale_divergence.png}
\end{minipage}\hfill
\begin{minipage}[t]{0.48\textwidth}
\centering
\includegraphics[width=\textwidth]{figures/2_curves.png}
\end{minipage}
\caption{One frozen representation, three objectives.
\emph{Left:} critic scale $\E|Q|$ (log axis) while training on a frozen goal embedding $\psi$.
The density objective (orange) grows without bound once the orthonormality term is removed: its linear term always rewards a larger matched score, and play data supplies the same matched pairs over and over.
The indicator-reward objective of \S\ref{sec:bounded} (aqua) converges to the Monte-Carlo estimate of the true discounted occupancy (dashed), because every entry of its TD target is a finite number.
\emph{Right:} task success over training under the same frozen embedding: the goal-reaching TD (blue) and the occupancy TD on the frozen phase-1 embedding (aqua) hold the goal-conditioned IQL level for the full run; the density TD (orange) degrades immediately.}
\label{fig:diverge}
\end{figure}

\section{Experiments}
\label{sec:experiments}

We organize the evaluation around four questions: does region-based TD restore goal-conditioned control (\S\ref{sec:rq1}); does occupancy fine-tuning improve held-out reward transfer (\S\ref{sec:rq2}); does the learned function satisfy measure properties (\S\ref{sec:rq3}); and does the approach extend to multi-object manipulation (\S\ref{sec:rq4}).

\paragraph{Setting.}
We use OGBench state-based manipulation \citep{park2024ogbench}, primarily \texttt{cube-single-play-v0}: a $6$-DoF arm, one cube, $28$-dimensional states, and $10^6$ offline transitions of scripted play that repeatedly picks the cube up, carries it, and sets it down.
No benchmark reward labels are used to fit any of our objectives.
The method is given a state abstraction $\cube(x)\in\mathbb{R}^3$ (cube position) and a radius $\epsilon=4$\,cm, which define goal regions $B_\epsilon(g)=\{x:\|\cube(x)-\cube(g)\|\le\epsilon\}$.
This is a substantive assumption, three of twenty-eight state dimensions, at the benchmark's own tolerance, and we quantify what it contributes in \S\ref{sec:results}.

\paragraph{Protocol.}
Unless stated otherwise: rank-$50$ bilinear critic $F^\top\psi$, twin heads with a minimum target, $\gamma=0.99$, hindsight-relabeled batch goals (a batch of $n$ transitions yields an $n\times n$ matrix of state-goal pairs), a flow-matching behavior-cloning actor with one-step distillation, $500$k gradient steps, $2$-$3$ seeds per cell.
\emph{Task success} means the benchmark's goal-conditioned evaluation, averaged over its five tasks with $50$ episodes each.
Measure-side evaluations (Monte-Carlo occupancy comparisons, reward-transfer protocols) are defined where they are used.

\subsection{Does region-based TD restore goal-conditioned control?}
\label{sec:rq1}\label{sec:results}

Yes.
Table~\ref{tab:main} compares goal-reaching success on the three environments, against the OGBench reference implementations and against baselines re-implemented in the matched architecture.

\begin{table}[t]
\centering\small
\begin{tabular}{l ccc}
\toprule
Method & \texttt{cube-s} & \texttt{cube-d} & \texttt{scene} \\
\midrule
\multicolumn{4}{l}{\emph{reported by OGBench \citep{park2024ogbench}, reference implementations, $8$ seeds}} \\
GCBC   & $0.06$ & $0.01$ & $0.05$ \\
GCIVL  & $0.53$ & $0.36$ & $0.42$ \\
GCIQL  & $0.68$ & $0.40$ & $0.51$ \\
CRL    & $0.19$ & $0.10$ & $0.19$ \\
HIQL   & $0.15$ & $0.06$ & $0.38$ \\
\midrule
\multicolumn{4}{l}{\emph{architecture-matched comparison, same data, flow-BC actor, bilinear rank-$50$ critic, $2$-$10$ seeds}} \\
GCIQL baseline           & $0.928$ & $0.265$ & $0.281$ \\
CRL                      & $0.311$ & ,     & ,     \\
Tuned FB                 & $0.42$  & ,     & $0.047$ \\
FB (reference reimpl.)   & $0.532$ & ,     & ,     \\
FB loss on regions       & $0.284$ & ,     & ,     \\
\midrule
Reachability value (phase 1) & $\mathbf{0.967}$ & $\mathbf{0.549}$ & $\mathbf{0.484}$ \\
Region occupancy (phase 2)   & $0.940$ & $0.358$ & $0.342$ \\
\bottomrule
\end{tabular}
\caption{Goal-reaching success ($500$k steps; five tasks, $50$ episodes each).
The degenerate \texttt{puzzle-3x3}, where the matched implementation solves only task 1, is omitted.
The top block quotes the OGBench paper's reference implementations (their architectures and tuning); the middle block re-implements baselines in the matched architecture, except FB (reference reimpl.), a reimplementation of the reference FB code (single seed).
Dashes: not run.}
\label{tab:main}
\end{table}

\paragraph{Region-based TD restores goal-conditioned control.}
The reachability value outperforms the architecture-matched GCIQL baseline on all three environments, with the largest improvements on \texttt{cube-double} and \texttt{scene} (Table~\ref{tab:main}).
Against OGBench's reference implementations it is the strongest evaluated method on \texttt{cube-single} and \texttt{cube-double} and comparable to reference GCIQL on \texttt{scene}.
With a bounded target, the same representation class that FB uses is competitive with the strongest evaluated baselines.

\paragraph{Successor-style baselines degrade where occupancy concentrates.}
CRL, the contrastive counterpart, scores $0.10$-$0.31$ across implementations, and FB drops to $0.047$ on \texttt{scene}, the environment where the object-at-rest concentration is strongest (\S\ref{sec:playdata}).
This pattern is consistent with the proposed mechanism, although it does not isolate it.

\paragraph{Which change is responsible?}
A $2\times2$ ablation crosses the loss form (FB's surrogate vs.\ indicator-TD) with the query support (point vs.\ region).
FB's own loss on regions reaches $0.284$ (Table~\ref{tab:main}, ``FB loss on regions''); indicator-TD with the region collapsed back to a point ($h=I$ on the batch) reaches $0.108$; the combination reaches $0.922$ (replicated in a second implementation).
Neither change suffices alone.
The region also need not be the evaluation variable: the same ball built on the \emph{end-effector} position, correlated with, but distinct from, the evaluated cube, reaches $0.823$.
The region variable therefore need not exactly match the evaluation variable, although this ablation does not establish that the choice of variable is unimportant in general.

\subsection{Does occupancy fine-tuning improve held-out reward transfer?}
\label{sec:rq2}\label{sec:reward-transfer}

Yes: fine-tuning nearly doubles compositional-reward return, and with the two corrections below the final model reaches $1.9\times$ the mean return of the reimplemented FB baseline (Table~\ref{tab:transfer}).

\begin{table}[t]
\centering\small
\begin{tabular}{l cccc}
\toprule
Method & concentrated & broad & compositional & mean \\
\midrule
FB (reference reimpl.)  & $8.4$ & $35.4$ & $6.9$ & $17.2$ \\
FB loss on regions      & $6.4$ & $27.2$ & $5.0$ & $13.1$ \\
\midrule
Reachability value (phase 1) & $12.1$ & $27.8$ & $8.5$  & $16.8$ \\
$+$ occupancy (phase 2)      & $16.3$ & $28.5$ & $15.8$ & $20.3$ \\
$+$ reward-embedding coverage & $20.2$ & $31.5$ & $23.0$ & $24.5$ \\
$+$ measure regularizers & $\mathbf{26.1}$ & $\mathbf{41.6}$ & $\mathbf{33.0}$ & $\mathbf{32.4}$ \\
\bottomrule
\end{tabular}
\caption{Held-out reward return on \texttt{cube-single}: $48$ rewards over the cube abstraction, absolute discounted return (rewards in $[0,1]$, returns in $[0,86.6]$, behavior-policy baseline $8.1$; family means of two seeds, FB single-seed).
Concentrated $=$ balls at radii $0.4/0.8$ and Gaussians at bandwidth $0.4$; broad $=$ ball $1.6$ and Gaussian $0.8$; compositional $=$ two disjoint zones.
Rows below the rule are cumulative; goal-reaching success is unchanged by the two corrections ($0.940\to0.947$-$0.950$).}
\label{tab:transfer}
\end{table}

The rewards are functions of the supplied cube abstraction, all with values in $[0,1]$: indicator balls at radii $0.4/0.8/1.6$ (the training kernel is $0.4$), Gaussians at bandwidths $0.4/0.8$, and a two-zone reward (two disjoint balls more than $1.6$ apart, not expressible as any single goal); eight instances per family, centers drawn from dataset states.
Task vectors come from least-squares reward inference over $10^4$ relabeled dataset rows, no ridge term, no candidate search.
For per-family analysis we score control as the fraction of the gap between a random-$z$ floor and a single-goal oracle proxy; scores above $1$ mean beating that proxy, not optimality.

\paragraph{FB's return comes from broad rewards.}
The reference-reimplementation row shows FB performing best on the broad family, above our phase-1 and phase-2 models, because a $16$\,cm ball assigns reward to almost any cube displacement, and $3$-$5\times$ below the final method on the concentrated and compositional families, where the reward has structure.

\paragraph{Occupancy fine-tuning trades goal-reaching success for transfer.}
Fine-tuning reduces goal-reaching success on \texttt{cube-double} and \texttt{scene} ($-0.19$/$-0.14$; \S\ref{sec:rq4} removes this reduction) but improves reward transfer where occupancy matters: on the compositional family, which no single goal expresses, return nearly doubles.
The two corrections below raise mean return further without reducing goal-reaching success.

\paragraph{Verifying the readout weight.}
The inverse-mass weight of Equation~\ref{eq:readout} is visible in the data.
With the weight omitted, predicted returns rank well in every family (Spearman $0.62$-$0.90$) but are uniformly $0.014$-$0.030$ of the Monte-Carlo magnitude, a single family-independent scale factor.
The mean region mass, measured independently, is $\bar N=0.0137$; the uncorrected predictions lie on $y=\bar N x$ (Figure~\ref{fig:scale}), and weighting the readout by $1/\Nhat(g)$ restores scale to $0.79$-$2.2$ (derivation in Appendix~\ref{app:derivations}).
The correction substantially reduces the scale error, although predictions can still differ from Monte Carlo by up to a factor of two.

\begin{figure}[t]
\centering
\includegraphics[width=0.46\textwidth]{figures/scale.pdf}
\caption{Predicted vs.\ Monte-Carlo return on all $48$ held-out rewards, log-log.
Uncorrected reward inference tracks $y=\bar Nx$ with $\bar N=0.0137$ the independently measured mean region mass; the $1/\Nhat$-weighted readout tracks $y=x$.}
\label{fig:scale}
\end{figure}

\paragraph{Covering reward-induced embeddings.}
Correcting the scale does not by itself improve behavior, two-zone control \emph{falls} from $0.86$ to $0.46$ under the corrected embeddings.
This is the training-side consequence anticipated in \S\ref{sec:readout}: an inferred embedding for a broad reward is an average of many goal embeddings and lies in a part of task space that training, which samples $z$ from random directions and single-goal embeddings, never visits.
The test is to add such embeddings to training: a third task-sampling component draws a random synthetic reward and embeds it as inference would.
A \textsc{kernel} generator matches the evaluation families' functional forms; an \textsc{rff} generator shares no functional form with any evaluation family and controls for reward-family leakage (Appendix~\ref{app:rewards}).

\begin{table}[t]
\centering\small
\begin{tabular}{l c c c c}
\toprule
Family & baseline & \textsc{kernel} $\eta{=}0.25$ & \textsc{kernel} $\eta{=}0.5$ & \textsc{rff} $\eta{=}0.25$ \\
\midrule
\textsc{ball} $0.4$   & $0.80$ & $0.87$ & $0.87$ & $0.79$ \\
\textsc{ball} $0.8$   & $0.66$ & $0.71$ & $\mathbf{0.95}$ & $0.73$ \\
\textsc{twozone}      & $0.86$ & $0.92$ & $\mathbf{1.35}$ & $1.11$ \\
\textsc{smooth} $0.4$ & $0.50$ & $0.49$ & $\mathbf{0.78}$ & $0.55$ \\
\textsc{smooth} $0.8$ & $0.15$ & $0.27$ & $0.33$ & $0.12$ \\
\textsc{ball} $1.6$   & $0.07$ & $\mathbf{0.53}$ & $\mathbf{0.49}$ & $\mathbf{0.55}$ \\
\midrule
mean (all 6)          & $0.507$ & $0.630$ & $\mathbf{0.795}$ & $0.640$ \\
task success          & $0.887$ & $0.948$ & $0.950$ & $0.930$ \\
set-mass calibration check (hw $0.8$) & $1.10/1.16$ & $1.10/1.19$ & $1.23/1.31$ & $1.26/1.12$ \\
\bottomrule
\end{tabular}
\caption{Zero-shot control (normalized return relative to random-$z$ and single-goal-oracle baselines) after adding reward-derived task embeddings to training at mixing fraction $\eta$; two-seed means.
The bottom rows check that the intervention reduces neither task success nor calibration.
Oracle ceilings agree within $3$-$23\%$ across these runs, so normalized scores are comparable among them.}
\label{tab:rewardz}
\end{table}

Training on reward-induced task embeddings improves control across most reward families (Table~\ref{tab:rewardz}).
The broad indicator ball, at $0.07$ at baseline, recovers to $0.49$-$0.55$ in every arm, including \textsc{rff}, whose training rewards share no form with the evaluation family: the improvement comes from populating the averaged region of task space, not from familiarity with the test rewards.
The compositional family reaches $1.35$, exceeding the better single-zone proxy by $35\%$, which is what reward composition should produce.
Combined with the regularizers of \S\ref{sec:rq3}, the broad Gaussian improves from $0.24/0.43$ to $0.46$-$0.86$ across seeds, though it remains the least reliable family.
Task success improves alongside ($0.887\to0.950$) and calibration is intact.
Figure~\ref{fig:bandwidth} resolves the same result by the reward's spatial support.

\begin{figure}[t]
\centering
\includegraphics[width=0.72\textwidth]{figures/bandwidth.pdf}
\caption{Zero-shot control against the reward's spatial support (normalized return as in Table~\ref{tab:rewardz}).
Control degrades as support widens within each family.
Adding reward-induced embeddings to the task-sampling distribution recovers the wide indicator ball and the compositional family, but not the wide Gaussian; the held-out reward form (\textsc{rff}) performs comparably to the matched one on the recovered families.}
\label{fig:bandwidth}
\end{figure}

\paragraph{What occupancy fine-tuning adds over the reachability value.}
\label{sec:cross-method}

Table~\ref{tab:transfer} answers on returns; the mechanism is visible when both policies are evaluated with the goal supplied directly: the phase-2 policy accumulates $1.13$-$2.00\times$ more reward than the phase-1 policy on identical rewards.
A reachability value is saturated once the region is entered, so its greedy policy enters and stops; a policy trained on occupancy returns to and remains in the rewarded region, exactly the post-arrival occupancy that \S\ref{sec:playdata} showed dominates this data.
The advantage disappears on the broadest rewards ($0.99$-$1.05\times$), where remaining in place earns nothing extra.
(The phase-2 checkpoint continues training from the phase-1 one, so the comparison is not iso-budget; the goal-supplied comparison is less exposed to this, but an equal-training-budget control has not been run.)

\subsection{Does the learned function satisfy measure properties?}
\label{sec:rq3}\label{sec:measure-validity}

Only partially: the estimator is coarsely calibrated but carries substantial negative mass, and removing the negative mass requires both regularizers of this section (Table~\ref{tab:measure}).
Strong task success could not have told us this: control depends mainly on the critic's action ranking at the current state.
Empirically, the goal-reaching and occupancy critics select the same greedy action on only $29.5\%/29.9\%$ of a grid of $24$ states $\times$ $150$ goals $\times$ $16$ candidate actions (chance $6.25\%$).
Task success provides evidence that the reachability critic induces an effective policy; it does not establish calibration of $M^{\Visit}$.
We therefore test the occupancy estimator directly for the properties a measure must have: calibrated set mass and non-negativity.

\subsubsection{Set-mass calibration}
\label{sec:set-calibration}

We compare predicted occupancy of random boxes against Monte-Carlo rollouts of the same policy ($32$ starts $\times$ $3$ rollouts, horizon $200$).
The learned quantity is a region mass, so converting it to a density requires dividing by the region's probability under the data distribution, $N(g)=\rho(B_\epsilon(g))$; a small trained head $\Nhat(g)$ estimates this factor.

\begin{table}[t]
\centering\small
\begin{tabular}{c c c c}
\toprule
box half-width & model/MC (exact $N$) & model/MC (trained $\Nhat$) & Spearman \\
\midrule
$0.4$ & $3.22/2.03$ & $1.94/1.22$ & $+0.31/+0.33$ \\
$0.8$ & $1.92/2.08$ & $\mathbf{1.10/1.16}$ & $+0.68/+0.72$ \\
$1.6$ & $2.33/2.18$ & $\mathbf{1.23/1.14}$ & $+0.85/+0.82$ \\
\bottomrule
\end{tabular}
\caption{Calibration per seed: predicted over Monte-Carlo set mass (ratio; $1.0$ is perfect) and rank correlation across boxes.
Monte Carlo truncates at $H{=}200$, omitting $\gamma^{H}=0.134$ of the discounted tail, so ratios near $1.15$ are near-exact.
Coarse boxes are calibrated; the smallest scale is not.}
\label{tab:calib}
\end{table}

At half-widths $0.8$ and $1.6$ the corrected estimate is within $10$-$23\%$ of Monte Carlo with rank correlation $0.68$-$0.85$ (Table~\ref{tab:calib}).
Two qualifications: with exact counts in the denominator, the uncorrected estimate is $1.9$-$3.2\times$ too large (the head's $+28\%$ bias partially compensates), and the smallest scale, the training kernel's own, remains $1.2$-$1.9\times$ off on every variant we trained, consistent with a resolution limit of the rank-$50$ factorization (\S\ref{sec:limitations}).

\subsubsection{Non-negativity and mass conservation}
\label{sec:nonneg}

A measure is non-negative.
The learned occupancy estimator is not: across $32$ policies and $20$k dataset goals, about half of all values are negative, the negative values carry $24$-$29\%$ of the total absolute mass at roughly a third of the positive values' magnitude, and the violation is \emph{largest} in the densest parts of goal space ($0.55$-$0.59$ of values negative in the top density quartile vs.\ $0.45$ in the bottom).
These are not boundary fluctuations; they are systematic oscillations where the regression is most constrained.
Calibration and invalidity coexist because signs cancel inside a box.

Replacing the discontinuous $0/1$ reward with a smooth kernel of matched mass changes nothing (negative fraction $0.50$-$0.52$; Appendix~\ref{app:validity}), so the violation is slack in the regression, not an artifact of the target's sharpness.

The direct correction is a one-sided penalty $\lambda\,\E[\mathrm{relu}(-M)^2]$.
Alone, it fails, and the failure identifies the missing constraint: at $\lambda\in\{1,10\}$ the negative fraction drops to $3.1\%/0.7\%$, but total mass inflates to $2.5$-$3.6\times$ Monte Carlo.
The penalty raises negative values, and no term in the objective constrains the total.
The missing constraint is a property of the region occupancy's fixed point: because the cube is always \emph{somewhere}, averaging region mass over goals, each weighted by the inverse of its region's data probability, counts total discounted time,
\begin{equation}
\label{eq:mass-anchor}
\E_{g\sim\rho}\!\left[\frac{M(s,a,g)}{N(g)}\right]=\sum_{t\ge0}\gamma^t =\frac{1}{1-\gamma} \quad\Longrightarrow\quad \mathcal{L}_{\mathrm{mass}} = \Big((1-\gamma)\,\tfrac{1}{n}\sum\nolimits_j M_{ij}/\Nhat_j-1\Big)^{2}.
\end{equation}
The penalty holds the model to this identity (exactness conditions in Appendix~\ref{app:derivations}).
With both regularizers, sign and scale are fixed jointly (Table~\ref{tab:measure}) and task success is unchanged ($0.94$-$0.95$ vs.\ control $0.94$-$0.96$).

\begin{table}[t]
\centering\small
\begin{tabular}{l ccc}
\toprule
Model & negative-mass share & set-mass ratio (hw $0.8$/$1.6$) & total mass \\
\midrule
unregularized               & $0.24$-$0.29$ & $1.10$-$1.23$ & unconstrained \\
non-negativity only        & $<0.01$        & ,             & $2.5$-$3.6\times$ \\
non-negativity $+$ mass    & $0.010$-$0.014$ & $1.15$-$1.29$ & $1.002$ \\
\bottomrule
\end{tabular}
\caption{Measure diagnostics on \texttt{cube-single} (two seeds).
Negative mass: share of total absolute mass carried by negative values.
Set-mass ratio: predicted over Monte-Carlo region mass with the trained $\Nhat$ (Table~\ref{tab:calib} resolves it by scale).
Total mass: training-time total-occupancy ratio (target $1$).
The non-negativity penalty alone removes sign errors but inflates total mass; adding the mass-conservation penalty fixes both.}
\label{tab:measure}
\end{table}
The measured cost: rank correlation at the middle calibration scale drops from $0.68$ to $\approx0.50$ (the coarse scale is unaffected); penalty-only runs already showed this degradation, implicating the penalty strength; a weaker penalty is under evaluation (\S\ref{sec:limitations}).

\subsection{Does the approach extend to multi-object manipulation?}
\label{sec:rq4}\label{sec:multienv-zeroshot}

Yes, with one change to how goals are sampled for the occupancy TD; the same change removes phase 2's reduction in goal-reaching success ($-0.19/-0.14$, Table~\ref{tab:main}).
On \texttt{cube-double} the trained region is a conjunction (both cubes within $\epsilon$), so a goal sampled uniformly from the dataset gives the occupancy TD a nonzero indicator at rate $0.046\%$, $28\times$ sparser than \texttt{cube-single}.
We tested the other candidate causes, kernel width, factorization rank, freezing the actor through the phase switch, and their compositions; eight pre-specified ablations, none of which improved success, and one change did: sampling the occupancy TD's goals by hindsight relabeling (future states of the same trajectory), which makes the indicator nonzero at the state-goal pairs the actor loss uses, by construction.
The cost: without an importance correction the hindsight-sampled object is no longer the successor measure with respect to the dataset distribution $\rho$ (\S\ref{sec:set-calibration}), and on \texttt{cube-single}, where dataset-sampled goals already give a nonzero indicator at $1.3\%$, the change reduces success by $0.17$-$0.24$.
It is an environment-conditional modification, not a default.

Table~\ref{tab:multienv-zeroshot} scores zero-shot control on the multi-object environments with the \texttt{cube-single} reward-transfer evaluation protocol extended to conjunctive regions (region centers drawn from visited states; the extended protocol reproduces the \texttt{cube-single} numbers of \S\ref{sec:reward-transfer} exactly).
Protocol as before: least-squares reward inference over relabeled rows, nearest-neighbor decoding of $z_r$ to a dataset goal state, the benchmark's success criteria, $25$ episodes per task at $500$k steps.

\begin{table}[t]
\centering\small
\begin{tabular}{l l cc cc}
\toprule
& & \multicolumn{2}{c}{benchmark tasks $\uparrow$} & \multicolumn{2}{c}{region reward $\uparrow$} \\
\cmidrule(lr){3-4}\cmidrule(lr){5-6}
Env & checkpoint & zero-shot & direct GC & final & anytime \\
\midrule
\texttt{cube-s} & reachability value & $0.968$ & $0.967$ & \multicolumn{2}{c}{$0.741$ (R1-R3 mean)} \\
\midrule
\texttt{cube-d} & reachability value & $0.480$ & $0.504$ & $0.15$ & $0.23$ \\
& $+$ occupancy, dataset goals     & $0.312$ & $0.368$ & $0.04$ & $0.04$ \\
& $+$ occupancy, hindsight goals & $0.488$ & $0.500$ & $0.08$ & $0.18$ \\
\midrule
\texttt{scene} & reachability value & $0.520$ & $0.468$ & $0.26$ & $0.42$ \\
& $+$ occupancy, dataset goals     & $0.376$ & $0.332$ & $0.21$ & $0.35$ \\
& $+$ occupancy, hindsight goals & $\mathbf{0.872}$ & $0.916$ & $\mathbf{0.45}$ & $\mathbf{0.74}$ \\
\bottomrule
\end{tabular}
\caption{Zero-shot control beyond \texttt{cube-single}.
``Zero-shot'' infers the task from reward-labeled dataset rows; ``direct GC'' is the benchmark's goal-conditioned evaluation of the same checkpoint.
Region rewards are conjunctive balls at the trained radii around visited-state centers, scored at the final step and at any step; the \texttt{cube-single} region row uses the single-object region protocol and is not comparable to the multi-object cells.
Single-seed results; the hindsight direct-GC numbers replicate at three seeds ($0.531$ \texttt{cube-d}, $0.927$ \texttt{scene}, Table~\ref{tab:main}).}
\label{tab:multienv-zeroshot}
\end{table}

\section{Related work}
\label{sec:related}

\paragraph{Successor representations and measures.}
Successor representations predict discounted future occupancy \citep{dayan1993improving}; successor measures generalize them to continuous spaces as measures, or densities against a reference distribution \citep{blier2021learning}; FB representations factorize that density and infer tasks by least-squares reward regression \citep{touati2021learning,touati2023does}.
We keep FB's bilinear form and zero-shot protocol and change the estimand from a pointwise density, whose empirical Bellman target is a Dirac measure, to region occupancy, whose target is a sparse indicator reward, a change whose necessity is specific to data that concentrates occupancy at resting states (\S\ref{sec:playdata}).

\paragraph{Goal-conditioned RL.}
Goal-reaching values with sparse indicator rewards and hindsight relabeling are the standard objects of offline GCRL \citep{kaelbling1993learning,andrychowicz2017hindsight,kostrikov2021offline,park2023hiql}.
Contrastive RL \citep{eysenbach2022contrastive} learns classifiers of future states rather than densities, avoiding the same singular target from a different direction; our diagnosis in \S\ref{sec:playdata} makes the connection explicit, FB's surrogate is contrastive in structure but lacks the saturating positive term that keeps classification bounded.
A goal-conditioned value evaluates each goal under that goal's own policy; a successor measure retains an independent query and must support reward composition, which is why our evaluation separates reachability (task success) from the measure claim (calibration, sign, transfer).

\paragraph{Hitting values versus occupancy.}
The absorbing and non-absorbing TD targets answer different questions, whether a region is reached versus how long it is occupied.
Manipulation benchmarks score first arrival into a tolerance region, so a reachability value suffices for control; Table~\ref{tab:main} shows occupancy is nonetheless the object that improves reward transfer.

\section{Limitations and open problems}
\label{sec:limitations}

\paragraph{The region is supplied.}
$\cube$ and $\epsilon$ come from the benchmark's success criterion.
The $2\times2$ ablation bounds what this buys, the region alone rescues nothing ($0.284$ under FB's loss), and a correlated non-evaluated region retains most of the benefit ($0.823$), but discovering outcome variables and scales from data is unsolved here.
All $48$ transfer rewards are functions of the supplied abstraction; ``zero-shot'' refers to held-out rewards, not held-out abstractions.

\paragraph{Multi-object evidence covers success, not calibration.}
Zero-shot control now has multi-object evidence (\S\ref{sec:multienv-zeroshot}), and phase 2's success reduction on \texttt{cube-double}/\texttt{scene} is repaired by hindsight goal sampling, but the occupancy-calibration evidence is weaker there: set-mass ratios are $1.57/1.14$ with rank correlations $0.37/0.50$ and $71$-$78\%$ of probe centers degenerate, and return calibration holds only at $2$-$4\times$ the trained radius.
The hindsight-sampled object that restores control is not the successor measure with respect to $\rho$; whether one object can be both a valid measure and a critic for policy improvement on these environments, by mixing query distributions, or by giving the actor its own hindsight-trained critic, is the open question.

\paragraph{Small-scale calibration, and the cost of the regularizers.}
Set mass at the training kernel's own scale is $1.2$-$1.9\times$ off on every variant, unchanged by kernel smoothing or penalties, consistent with a resolution limit of the rank-$50$ factorization; raising the rank is the natural next experiment.
The regularizers cost mid-scale rank correlation ($0.68\to0.50$); a weaker penalty is under evaluation.

\paragraph{Comparisons.}
The OGBench block of Table~\ref{tab:main} is quoted from the paper's reference implementations ($8$ seeds, their architectures and tuning) and is not an apples-to-apples ablation of our stack, notably, reference GCIQL solves \texttt{puzzle-3x3} ($0.95$) where our matched bilinear-flow stack does not.
None of the baselines has been run through the reward-transfer evaluation (Table~\ref{tab:main} dashes), so the return columns compare only our objectives against FB's loss at matched support.
The phase-2-vs-phase-1 return comparison does not use matched training budgets; the oracle proxy understates ceilings for broad and multimodal rewards; the random-$z$ floor is a single draw.

\paragraph{Statistical power; scope of theory; one actor class.}
Most cells have two or three seeds.
The contraction result covers fixed-policy operators, not the full system with a changing actor and twin-min targets.
All primary cells share one flow-matching BC actor; a single-seed sweep with a different actor reversed one comparison at one setting, so actor isolation remains open.

\section{Conclusion}
\label{sec:conclusion}

Manipulation play data has a structure that the domains where zero-shot RL methods were developed do not: the manipulated object is at rest half the time, so most of the discounted occupancy accrues while the object rests and carries no control signal.
A pointwise density objective has no finite regression target exactly there, and the standard surrogate fails on a frozen, known-good representation with and without its orthonormality term.
Matching the algorithm to the data resolves the failure: region queries make the objective ordinary goal-conditioned TD, and a two-stage procedure trains the reachability value, which carries the control signal, before adding the post-arrival occupancy that reward integration needs.
Evaluated jointly on goal-reaching and reward transfer (Table~\ref{tab:main}), the result beats the matched goal-conditioned baseline on every environment while earning $1.9\times$ FB's return on held-out rewards, with two further corrections, an inverse-mass weight for return prediction and reward-derived task vectors for policy coverage, accounting for most of the remaining transfer gap.
The remaining supplied input is the region itself, three coordinates and a radius, and learning it from data is the problem to solve next.

\bibliography{references}
\bibliographystyle{iclr2027_conference}

\appendix
\section{Additional experimental detail}
\label{app:detail}

\subsection{Play-data statistics}
\label{app:playstats}

Computed on \texttt{cube-single-play-v0} ($10^6$ transitions), cube slice of the raw observation.
Stationarity: $\|\Delta\cube\|<10^{-6}$ on $47.9\%$ of within-episode transitions ($51\%$ below $10^{-3}$); stationary runs have median length $44$ steps (mean $43.9$, p90 $47$).
Region mass: a $4$\,cm ball around a random dataset state contains on average $1.34\%$ of the data ($\bar N=0.0137$).
Post-arrival occupancy: discounted occupancy of the $4$\,cm goal ball accumulated after first arrival has dataset-level mean $\approx24.7$ (within-trajectory estimate $\approx30.6$, bimodal: passed-through vs.\ settled).
A goal-only regressor of this mass fit on the frozen $\psi(g)$ predicts Monte-Carlo occupancy at Spearman $0.60$-$0.63$ uniformly across first-arrival-time buckets; the reachability value alone predicts at $0.06$.

Multi-object play datasets ($200$ episodes sampled per dataset, conjunctive regions from the environments' success criteria).
\texttt{cube-double-play-v0}: each cube stationary on $72\%$ of transitions (all objects simultaneously on $48\%$); conjunctive region mass $3.3\times10^{-4}$; post-arrival occupancy of the conjunctive region mean $32.2$ (median $43$).
\texttt{scene-play-v0}: cube stationary $74\%$, drawer $91\%$, window $91\%$, buttons $99.6\%$ (all objects $60\%$); conjunctive region mass $3.7\times10^{-3}$; post-arrival occupancy mean $37.0$ (median $37.6$).

\subsection{Divergence: the full control set}
\label{app:diverge}

Section~\ref{sec:playdata} reports the frozen-representation experiment in condensed form.
The complete set of controls:

\begin{itemize}
\item \textbf{Frozen $\psi$, density surrogate.} Critic scale $154\to5.4\times10^{4}$ within $65$k steps; mean success decreases monotonically from $0.98$ (initial) through $0.60$ to $0.42$ by $100$k steps (three seeds).
The representation comes from a checkpoint supporting $0.97$ control and is held fixed, so representation quality cannot account for the outcome.
\item \textbf{Trainable $\psi$, density surrogate.} Collapses faster: $0.97\to0.07$ and $0.98\to0.06$ by $50$k (two seeds).
\item \textbf{Frozen $\psi$, indicator-reward TD.} Scale converges to the Monte-Carlo estimate of true discounted occupancy; both seeds complete $500$k steps without instability.
\item \textbf{Orthonormality retained.} Scale bounded; success $0.225$.
\end{itemize}

\subsection{Qualitative rollout}

\begin{figure}[t]
\centering
\includegraphics[width=0.95\textwidth]{figures/strip_vis_task2_ok.png}
\caption{A successful rollout of the fine-tuned occupancy policy (\texttt{cube-single}, benchmark task 2): approach, grasp, transport, and release inside the goal region.}
\label{fig:rollout}
\end{figure}

\subsection{A tabular instance of the same structure}
\label{app:playpen}

The estimator claims can be exhibited exactly in a $50$-state MDP with the same anatomy as manipulation play data: a $7\times7$ grid plus a pit, where the column is the outcome variable (the region coordinate) and the row is an interaction mode invisible to the region.
A $2\times2$ sticky bowl models a cube at rest, the dwell structure of \S\ref{sec:playdata} in miniature; the top row is a conveyor drifting two columns per step regardless of action (a carried cube whose position cannot be held); riding off its end drops into a hard-to-leave pit.

Fitting a rank-$8$ bilinear $F^{\top}\psi$ on $20$k offline transitions with identical architecture and data, indicator-reward regression reaches relative error $0.05$-$0.09$ and per-state goal-ranking correlation $0.96$-$0.99$; the density-TD form reaches only $0.38$ and $0.65$, and adding the orthonormality regularizer makes it worse ($0.53$ and $0.45$).
With no actor in the loop the density estimator fits poorly but does not diverge, consistent with the divergence of \S\ref{sec:playdata} being a property of the estimator-plus-actor system.
The environment also exhibits the ranking disagreement that motivates the occupancy objective: with a stable and a transient region weighted to be equally reachable, ranking by reachability value selects the transient region while ranking by occupancy selects the stable one; exact policy evaluation confirms the latter ($96.8$ vs.\ $47.9$ discounted return).

\subsection{Non-negativity and mass-conservation diagnostics, in detail}
\label{app:validity}

Negativity measurement: $32$ draws of $(s_0,a_0,z)$ with $z$ the projected embedding of a random dataset state and $a_0$ the actor's own action; $20$k dataset goals per draw.
Reported: fraction of negative values, share of $\sum|M|$ carried by negative values, both split by quartile of the goal's kernel density $N(g)$ over $10^5$ reference rows.
Baseline: negative fraction $0.48$-$0.53$, negative mass share $0.24$-$0.29$, rising with density quartile.
Penalty only ($\lambda{=}1/10$): negative fraction $3.1\%/0.7\%$, box mass $2.5$-$3.6\times$ MC.
Penalty ($\lambda{=}1$) $+$ mass conservation (coefficient $1$): negative fraction $0.078$-$0.096$, negative mass share $0.010$-$0.014$, flat across quartiles; training total-mass ratio $1.002$; set-mass ratios (trained-$\Nhat$) $1.15$-$1.29$ at half-widths $0.8/1.6$; Spearman $0.47$-$0.51$ (hw $0.8$) and $0.83$-$0.86$ (hw $1.6$); task success $0.940/0.953$ vs.\ control $0.937/0.963$ under the same parser.
Kernel control: a $C^1$ biweight kernel $(1-(d/h)^2)^2$, $h=0.6$ (mass-matched to the $0.4$ ball), leaves the negative fraction at $0.50$-$0.52$ and does not improve small-box ratios.

\subsection{Reward construction}

\begin{figure}[t]
\centering
\includegraphics[width=0.95\textwidth]{figures/5_rewards.png}
\caption{The held-out reward families on the cube workspace (cube-position plane; $\epsilon$-balls to scale).
The first five panels are functions of cube position.
The \textsc{linear} panel colors dataset states by $\psi(x)^{\top}\tilde z$ for one random direction of the trained embedding: the field is spatially incoherent in cube space, which is the visual form of the coverage problem of \S\ref{sec:readout}.}
\label{fig:rewardfams}
\end{figure}

\label{app:rewards}

All transfer rewards are functions of the cube position $\cube(x)$, with centers drawn uniformly from dataset states; eight instances per family; all values in $[0,1]$.
\textsc{ball}~$\epsilon$ is $\ind{\|\cube(x)-c_0\|\le\epsilon}$ for $\epsilon\in\{0.4,0.8,1.6\}$ in observation units (the training kernel is $0.4$); \textsc{smooth}~$\sigma$ is $\exp(-\|\cube(x)-c_0\|^2/2\sigma^2)$ for $\sigma\in\{0.4,0.8\}$; \textsc{twozone} is $\max_k\ind{\|\cube(x)-c_k\|\le0.4}$ with $\|c_1-c_2\|>1.6$.
Inference uses $m=10{,}000$ relabeled rows.
Prediction is scored from $12$ teleported starts $\times$ $2$ rollouts; control over $10$ episodes per drive at horizon $200$, so returns lie in $[0,\sum_{t<200}\gamma^t]=[0,86.6]$.
Table~\ref{tab:main}'s column groups average family means: concentrated $=$ \{\textsc{ball} $0.4$, \textsc{ball} $0.8$, \textsc{smooth} $0.4$\}, broad $=$ \{\textsc{ball} $1.6$, \textsc{smooth} $0.8$\}, compositional $=$ \textsc{twozone}.
Synthetic training rewards (\S\ref{sec:reward-transfer}) are drawn per batch: \textsc{kernel} mixes balls, Gaussians, and two-zone rewards with centers from the batch and log-uniform radii in $[0.2,2.0]$; \textsc{rff} draws $r=\mathrm{relu}(\cos(w^{\top}\cube+\varphi))$ with $w\sim\mathcal{N}(0,I)/b$, $\varphi\sim\mathrm{U}[0,2\pi)$.
Both are normalized before reward inference.

\subsection{The two-stage procedure is not an exact factorization}
\label{app:twostage}

For a singleton goal, occupancy factors as first-arrival probability times a recurrence term.
For finite regions this identity fails: the region can be re-entered from many boundary states, so recurrence depends on where the region was entered and on the actions taken inside it.
Empirically, recurrence remains action-dependent on our data: across a grid of states, goals, and candidate actions, the intraclass correlation of recurrence across actions is $0.24$-$0.29$, far from the action-independence an exact factorization would need.
The two-stage procedure therefore uses phase 1 only to initialize phase 2, which optimizes occupancy directly with TD.

\subsection{Freezing one factor: relation to standard stabilization methods}
\label{app:freeze}

Holding part of a value function fixed while the rest trains appears in several standard forms.
Semi-gradient TD stops gradients through the bootstrap target \citep{sutton2018reinforcement}.
Deep value learning holds a target copy of the network fixed between synchronizations \citep{mnih2015human}, and DDPG trains its actor against a critic whose target copies move on a slower timescale \citep{lillicrap2016continuous}.
Two-timescale analyses of actor-critic make the requirement precise: convergence arguments hold when one component evolves slowly enough to appear stationary to the other \citep{borkar1997stochastic,konda2000actor}.
Successor-feature methods fix the feature map and train the other factor of a factored value function by TD \citep{barreto2017successor}.
Phase 2's freeze differs from all of these in one respect: the goal embedding is fixed for the entire phase rather than updated slowly, which is closest to the successor-feature setting.

\subsection{Derivations}
\label{app:derivations}

\paragraph{Proof of Proposition~\ref{prop:fixed-points}.}
\emph{No density:} a measure has a density with respect to $\rho$ only if it is absolutely continuous with respect to $\rho$.
The point mass $\delta_{f(s,a)}$ assigns mass $1$ to the set $\{f(s,a)\}$, which has $\rho$-measure zero when $\rho$ is non-atomic, so $\delta_{f(s,a)}\not\ll\rho$ and no Radon-Nikodym derivative exists.
\emph{Boundedness and contraction:} for measurable $G$ and deterministic dynamics, $P(G\mid s,a)=\ind{f(s,a)\in G}\in\{0,1\}$.
Both region operators act on bounded functions with the sup norm.
For the absorbing operator $(\mathcal{T}^{\Reach}Q)(s,a,g)=\E[h+\gamma(1-h)Q(s',a',g)]$ and any $Q_1,Q_2$,
\begin{equation*}
|\mathcal{T}^{\Reach}Q_1-\mathcal{T}^{\Reach}Q_2| \le\gamma\,\E[(1-h)\,|Q_1-Q_2|]\le\gamma\|Q_1-Q_2\|_\infty
\end{equation*}
since $0\le 1-h\le 1$; the non-absorbing operator is a $\gamma$-contraction by the same bound without the $(1-h)$ factor.
Banach's fixed-point theorem gives a unique fixed point for each.
For the non-absorbing operator, unrolling the recursion and using the Markov property gives
\begin{equation*}
M(s,a,g)=\E[\sum_{t\ge0}\gamma^{t}h(s_{t+1},g)],
\end{equation*}
bounded by $1/(1-\gamma)$.
For the absorbing operator, let $T_g=\min\{t\ge0:s_{t+1}\in B_\epsilon(g)\}$ and $Q^{*}=\E[\gamma^{T_g}\ind{T_g<\infty}]$.
If $s_{1}\in B_\epsilon(g)$ then $T_g=0$ and the right side of the backup contributes $h=1=\gamma^{0}$; otherwise the backup contributes $\gamma\,\E[\gamma^{T_g'}\ind{T_g'<\infty}]$ with $T_g'$ the hitting time from $s_{1}$, and $T_g=T_g'+1$ on that event.
So $Q^{*}$ satisfies the fixed point equation, and by uniqueness it is the fixed point.
\hfill$\square$

\paragraph{The total-occupancy identity and its exactness conditions.}
For any start $(s,a)$, exchanging the expectation over goals with the sum over time,
\begin{equation*}
\E_{g\sim\rho}\!\left[\frac{M^{\pi}(s,a,B_\epsilon(g))}{N(g)}\right] =\E_{\pi}\Big[\sum_{t\ge0}\gamma^{t}\,w(s_{t+1})\Big], \qquad w(x):=\int_{B_\epsilon(x)}\frac{\rho(dg)}{N(g)},
\end{equation*}
using the symmetry $\ind{x\in B_\epsilon(g)}=\ind{g\in B_\epsilon(x)}$.
If goal regions partition the space (disjoint cells), $w\equiv1$ and the identity $\E_{g}[M/N]=1/(1-\gamma)$ is exact.
For overlapping balls, $w(x)=1$ whenever $N$ is constant on $B_\epsilon(x)$, so the identity holds to first order when $N$ varies slowly at scale $\epsilon$; and $\E_{x\sim\rho}[w(x)]=\int\rho(dg)\,N(g)/N(g)=1$ exactly by Fubini, so the identity is exact on average over states drawn from $\rho$.
The training penalty enforces the batch-mean form, which sits between these regimes.

\paragraph{The $1/N$ readout weight.}
Any reward's value is $Q_{r}(s,a)=\E_{x\sim\rho}[r(x)\,m(s,a,x)]$ with $m=dM/d\rho$ wherever the density exists.
The trained head outputs region mass,
\begin{equation*}
F^{\top}\psi(g)\approx M(s,a,B_\epsilon(g))=\int_{B_\epsilon(g)}m\,d\rho \approx m(s,a,g)\,N(g)
\end{equation*}
for $m$ approximately constant on the ball.
Hence $m(s,a,g)\approx F^{\top}\psi(g)/N(g)$ and
\begin{equation*}
Q_{r}(s,a)\approx\frac{1}{n}\sum_{j}r(g_{j})\,F^{\top}\psi(g_{j})/N(g_{j})
\end{equation*}
over goals $g_{j}\sim\rho$.
Omitting the $1/N(g_{j})$ factor multiplies the estimate by approximately $\bar N=\E[N]$ when $N$ is uncorrelated with $r\cdot m$ across goals, the $y=\bar N x$ line of Figure~\ref{fig:scale}.

\subsection{Monte-Carlo evaluation protocol}
\label{app:mc}

Teleport-style Monte-Carlo evaluation of the kind used in \S\ref{sec:set-calibration} must reset the environment before every rollout, not merely overwrite the simulator state: the episode step counter otherwise persists across rollouts and every rollout after the first truncates immediately, inflating model-to-Monte-Carlo ratios by one to two orders of magnitude.
A useful invariant when debugging such a harness: the implied total Monte-Carlo mass must approach $\sum_k\gamma^k$, because the tracked object is always somewhere.

\end{document}